\chapter{Oltre ASLG-PC12: altri dataset e altri compiti}
\label{cap:altri-dataset}

I capitoli precedenti mostrano che su ASLG-PC12 la domanda ``il GRPO aiuta?'' ha
poco margine per una risposta: il corpus è sintetico e una regola lo satura
(sezione~\ref{sec:sintetico}). Il progetto, però, non riguarda il gloss in sé: è
la fase iniziale di una domanda più generale, cioè \textbf{se il GRPO possa
allineare un modello a produrre uscite strutturate che rispettano vincoli rigidi,
con o senza supporti} (esempi nel prompt, Trie in decodifica, SFT preliminare).

Per rispondere servono compiti in cui il margine esiste e in cui la struttura
dell'uscita è diversa. Questo capitolo porta la stessa pipeline (stesso modello,
stessi iperparametri GRPO, stesso disegno a celle) su tre dataset: un corpus di
gloss annotato da umani, PHOENIX-2014T, e i due compiti non-gloss di
GrammarRL~\cite{tuccio2026grammarrl}, la classificazione gerarchica WOS-46985 e
il riconoscimento di entità CoNLL-2003. GrammarRL valuta la decodifica vincolata
da grammatica con reinforcement learning proprio su questi tre compiti; qui i
dataset sono gli stessi ma i numeri non sono confrontabili uno a uno, perché
split, modello, linearizzazione del target e reward sono diversi.

\section{I tre compiti}

\begin{table}[h]
\centering
\footnotesize
\begin{tabular}{llllr}
\toprule
Dataset & Compito & Ingresso & Uscita & Train / val / test \\
\midrule
PHOENIX-2014T & traduzione in gloss & frase tedesca & gloss DGS & $7.096$ / $519$ / $642$ \\
CoNLL-2003    & entità nominate & frase inglese & \texttt{TIPO:entità} & $14.041$ / $3.250$ / $3.453$ \\
WOS-46985     & classificazione & abstract inglese & \texttt{DOMINIO Area} & $\approx 42.300$ / --- / $4.699$ \\
\bottomrule
\end{tabular}
\caption{I tre dataset aggiunti. PHOENIX e CoNLL usano gli split ufficiali;
WOS non ne ha, e si usa la stessa logica di ASLG-PC12 (90/10 con seme 42).}
\label{tab:altri-dataset}
\end{table}

Tutti e tre si scaricano da soli al primo uso, da copie su Hugging Face fissate
per revisione e verificate con sha256 contro gli originali. Si usa solo il testo:
di PHOENIX non si scaricano i video.

\subsection{PHOENIX-2014T: gloss annotati da umani}

PHOENIX-2014T~\cite{camgoz2018phoenix} raccoglie le previsioni del tempo della
televisione tedesca, con la trascrizione in tedesco e la trascrizione in gloss
della lingua dei segni tedesca (DGS) fatta da annotatori umani. Un esempio reale
del test set:

\begin{center}
\small
\begin{tabular}{ll}
\toprule
Tedesco & \texttt{am donnerstag regen in der nordhälfte in der südhälfte mal sonne} \\
        & \texttt{mal wolken ähnliches wetter dann auch am freitag} \\
Gloss   & \texttt{DONNERSTAG NORDWEST REGEN REGION SONNE WOLKE WECHSELHAFT} \\
        & \texttt{DANN FREITAG AEHNLICH WETTER} \\
\bottomrule
\end{tabular}
\end{center}

A differenza di ASLG-PC12 il gloss \emph{non} è l'ingresso maiuscolizzato: le
parole funzionali spariscono, l'ordine cambia, intere espressioni diventano un
segno (\emph{mal sonne mal wolken} $\to$ \texttt{SONNE WOLKE WECHSELHAFT}),
le dieresi diventano digrammi (\texttt{AEHNLICH}). La sovrapposizione
fra ingresso e uscita è del $15{,}59\%$ contro il $98\%$ di ASLG-PC12
(sezione~\ref{sec:sintetico}). È il corpus su cui la letteratura riporta i
migliori risultati intorno a ROUGE-L $55{,}18$~\cite{ouargani2023t2g}.

Rispetto ad ASLG cambia solo il profilo del prompt (tedesco $\to$ DGS); il
vocabolario chiuso del Trie conta circa $1.080$ gloss estratti dal train, e
tutto il resto (reward, Trie a doppia radice, metriche) è invariato.

\subsection{CoNLL-2003: riconoscimento di entità}

CoNLL-2003~\cite{tjongkimsang2003conll} è il riferimento classico per il
riconoscimento di entità nominate: frasi di agenzia Reuters in cui ogni entità è
marcata come persona (\texttt{PER}), organizzazione (\texttt{ORG}), luogo
(\texttt{LOC}) o altro (\texttt{MISC}). In GrammarRL il modello genera un JSON;
qui il target è \emph{linearizzato} in una sequenza di token \texttt{TIPO:entità},
con le parole dell'entità unite da trattini bassi, ordinate per tipo e senza
ripetizioni, oppure il solo \texttt{NONE}. Esempi reali del test set:

\begin{center}
\small
\begin{tabular}{ll}
\toprule
Frase & Target \\
\midrule
\texttt{8. Dominick Gauthier ( Canada ) 22.73} & \texttt{PER:Dominick\_Gauthier LOC:Canada} \\
\texttt{Japan then laid siege to the Syrian \dots} & \texttt{LOC:Japan MISC:Syrian} \\
\texttt{AL-AIN , United Arab Emirates} & \texttt{LOC:AL-AIN LOC:United\_Arab\_Emirates} \\
\texttt{on Friday :} & \texttt{NONE} \\
\bottomrule
\end{tabular}
\end{center}

Il $20\%$ delle frasi valutate non contiene entità. Il compito è
strutturalmente diverso dal gloss: l'uscita non traduce l'ingresso, ne
\emph{seleziona} dei pezzi e li etichetta.

\subsection{WOS-46985: classificazione gerarchica}

Web of Science~\cite{kowsari2017hdltex} raccoglie $46.985$ abstract scientifici
etichettati con un dominio (sette: informatica, ingegneria elettronica,
psicologia, ingegneria meccanica, ingegneria civile, medicina, biochimica) e
un'area di quel dominio, per $134$ coppie (dominio, area) in tutto. Il target è
linearizzato in due token: il dominio e l'area con le parole unite da trattini
bassi. Un esempio reale del test set:

\begin{center}
\small
\begin{tabular}{p{0.68\textwidth}l}
\toprule
Abstract (inizio) & Target \\
\midrule
\texttt{This paper is mainly discuss about software component testing as an
interesting challenge in Component Based Software Engineering (CBSE) area.
\dots} & \texttt{CS Software\_engineering} \\
\bottomrule
\end{tabular}
\end{center}

Le etichette sono quelle canoniche del codice ufficiale delle classi.
Gli abstract sono lunghi: vengono troncati alle prime $256$ parole, così il
prompt resta entro la lunghezza massima. Per lo stesso motivo il dataset si usa
\textbf{solo in zero-shot}: tre abstract di esempio nel prompt non ci
starebbero. Con due soli token, ROUGE-L e BLEU dicono poco: la metrica da
leggere è l'exact match (dominio e area entrambi giusti), e il gloss-F1 dà
credito parziale al solo dominio giusto.

\section{Un vincolo per ogni compito}
\label{sec:vincolo-per-compito}

Il Trie a doppia radice (sezione~\ref{sec:trie-dual-root}) ammette qualunque
sequenza di parole di un vocabolario chiuso. Va bene per il gloss, non per gli
altri due compiti, e per due ragioni opposte: su CoNLL è troppo stretto, su WOS
troppo largo. La grammatica dell'uscita diventa quindi una scelta del config
(\texttt{grammar.mode}), e la \emph{stessa} definizione guida il decoder, la
validity e la reward di formato, così ciò che il decoder impone e ciò che le
metriche chiamano ``valido'' non possono divergere.

\begin{itemize}
  \item \textbf{\texttt{vocab}} (ASLG-PC12, PHOENIX-2014T): il Trie storico sul
        vocabolario di train, invariato.
  \item \textbf{\texttt{source\_spans}} (CoNLL-2003): un Trie \emph{per frase},
        costruito dal prompt stesso. Ogni parola dell'uscita è
        \texttt{TIPO:pezzo}, dove il pezzo è una sequenza contigua di al più
        dieci parole della frase; in alternativa, il solo \texttt{NONE}. Con il
        vocabolario di train il $48\%$ delle entità del test sarebbe stato
        impossibile da generare (sono nomi propri mai visti); con il Trie per
        frase ogni entità gold è raggiungibile e non c'è alcuna fuga di
        informazione dal test.
  \item \textbf{\texttt{sequences}} (WOS-46985): l'uscita intera deve essere
        una delle $134$ coppie viste nel train. Il Trie a token libero lasciava
        generare \texttt{Medical Medical}, l'area prima del dominio, o un'area
        di un altro dominio; ora il decoder impone dominio, poi un'area di quel
        dominio, poi la fine.
\end{itemize}

\subsection{Ogni entità una sola volta}
\label{sec:span-unici}

La prima campagna su CoNLL ha mostrato che la grammatica \texttt{source\_spans}
lasciava aperto un varco, e che il GRPO lo ha trovato. Con il Trie attivo il
GRPO dal modello base ha imparato a scrivere la stessa entità con tutti i tipi:

\begin{center}
\small
\begin{tabular}{ll}
\toprule
Gold      & \texttt{PER:Oleg\_Shatskiku} \\
Generato  & \texttt{PER:Oleg\_Shatskiku ORG:Oleg\_Shatskiku LOC:Oleg\_Shatskiku} \\
\bottomrule
\end{tabular}
\end{center}

È \emph{reward hacking} nel senso stretto: se il modello non sa il tipo, coprirsi
con tutti e tre garantisce credito parziale da ROUGE-L e BLEU (qui ROUGE-L
$0{,}5$) invece di rischiare zero. Il $53\%$ delle uscite del GRPO zero-shot
ripeteva un'entità, contro il $25\%$ del modello base e meno dell'$1\%$
dell'SFT. Le metriche lo mostrano in modo netto: ROUGE-L sale da $0{,}27$ a
$0{,}35$, l'exact match \emph{si dimezza} da $0{,}045$ a $0{,}023$, e la validity
($0{,}998$) non se ne accorge, perché l'uscita era ammessa dalla grammatica.

Una prima correzione ha vietato il \emph{pezzo identico}: ogni pezzo della
frase una sola volta, con un solo tipo. Ha funzionato solo a metà. Nella
campagna rilanciata il GRPO ha trovato subito il varco accanto, scrivendo pezzi
\emph{annidati}:

\begin{center}
\small
\begin{tabular}{ll}
\toprule
Frase    & \texttt{Japan then laid siege to the Syrian penalty area \dots} \\
Gold     & \texttt{LOC:Japan MISC:Syrian} \\
Generato & \texttt{PER:siege PER:siege\_to PER:siege\_to\_the MISC:game} \\
\bottomrule
\end{tabular}
\end{center}

Con la regola dei pezzi identici, i pezzi annidati erano nel $21$--$24\%$ delle
uscite delle celle GRPO zero-shot contro l'$1{,}6\%$ dell'SFT, e le uscite
avevano in media $1{,}9$ volte i token del gold. Nella cella SFT$\to$GRPO,
prima la migliore, l'exact match crollava da $0{,}78$ a $0{,}09$: la validity
restava a $0{,}98$, perché la grammatica ammetteva tutto questo.

La correzione definitiva sta ancora nella grammatica, non nella reward:
\textbf{le entità non si sovrappongono}. Ogni parola della frase può stare in
più entità solo quante volte compare nella frase (il suo \emph{budget}). Il
gold rispetta sempre il budget, perché le sue entità sono porzioni disgiunte
della frase: su tutte le $20.744$ frasi di train, validation e test di
CoNLL ogni gold è ammesso (anche \texttt{PER:Costas\_Simitis
PER:Costas\_Skandalidis}, dove \emph{Costas} compare due volte nella frase),
mentre \texttt{PER:siege PER:siege\_to} è impossibile, perché \emph{siege}
compare una volta sola. Nel decoder ogni nodo del Trie per frase conosce i
pezzi che terminano nel suo sottoalbero; un ramo nessuno dei cui pezzi sta
ancora nel budget viene potato, e una parola può chiudersi solo se il suo pezzo
ci sta. Così il decoder non arriva mai in un vicolo cieco. Lo stesso vincolo
entra nella validity e nella reward di formato, quindi penalizza l'hack anche
nelle celle senza Trie.

I risultati di CoNLL riportati più sotto sono quelli della \emph{prima}
campagna, prima di qualunque correzione; la campagna con il vincolo definitivo
è in corso.
\todo{aggiornare la tabella~\ref{tab:risultati-conll} con la campagna rilanciata}

\section{Adattamenti della reward}

Per PHOENIX-2014T la reward è quella di ASLG-PC12 (Capitolo~\ref{cap:reward}).
Per CoNLL e WOS i target sono cortissimi (uno o due token: \texttt{NONE},
\texttt{PER:Mary}, \texttt{CS Machine\_learning}), e i pesi di ASLG non
funzionano: la componente di struttura gold vale $-1$ fisso sotto i due token, e
la componente scalata dal verificatore, che la moltiplica, pure. Una risposta
corretta di un token prendeva $+0{,}30$ invece di $+1$. Su questi due compiti
la plausibilità dei bigrammi non porta informazione, e i due pesi passano
all'ordine dei gloss (distanza di edit a token, sensibile al maiuscolo, l'unica
componente che premia solo l'etichetta esatta):

\begin{center}
\small
\begin{tabular}{lrr}
\toprule
Componente & ASLG / PHOENIX & CoNLL / WOS \\
\midrule
traduzione (ROUGE-L) & $0{,}20$ & $0{,}20$ \\
BLEU                 & $0{,}20$ & $0{,}20$ \\
struttura gold       & $0{,}20$ & $0$ \\
ordine dei gloss     & $0{,}10$ & $0{,}40$ \\
verificatore         & $0{,}10$ & $0$ \\
formato              & $0{,}10$ & $0{,}10$ \\
ripetizione          & $0{,}10$ & $0{,}10$ \\
\bottomrule
\end{tabular}
\end{center}

\section{Protocollo}

Le celle sono quelle di ASLG-PC12: modello base (zero-shot, few-shot,
zero-shot senza Trie), solo SFT, solo GRPO, SFT seguito da GRPO, GRPO senza Trie
(few-shot e zero-shot) e SFT con obiettivo di massa ammessa. WOS ha solo le
celle zero-shot. Ogni GRPO fa $5.000$ passi, cioè circa $0{,}7$ epoche di
PHOENIX, $0{,}36$ di CoNLL e $0{,}12$ di WOS; gli SFT si fermano per early
stopping. La valutazione genera cinque uscite campionate per prompt
(temperatura $0{,}7$) su tutto il test di PHOENIX ($642$ frasi) e su $2.000$
prompt del test di CoNLL e WOS. Un solo seme per cella, come per ASLG-PC12
(sezione~\ref{sec:determinismo}).

\section{Risultati}

\subsection{PHOENIX-2014T}

\begin{table}[h]
\centering
\small
\begin{tabular}{lrrrr}
\toprule
Cella & ROUGE-L & Exact match & Non-copy & Validity \\
\midrule
Base zero-shot                 & $0{,}0961$ & $0{,}0000$ & $0{,}1083$ & $0{,}9586$ \\
Base few-shot                  & $0{,}3103$ & $0{,}0106$ & $0{,}2923$ & $0{,}9956$ \\
Base zero-shot senza Trie      & $0{,}1528$ & $0{,}0000$ & $0{,}0000$ & $0{,}0000$ \\
SFT                            & $0{,}4854$ & $0{,}0184$ & $0{,}3812$ & $0{,}9969$ \\
GRPO zero-shot                 & $0{,}2030$ & $0{,}0003$ & $0{,}1658$ & $0{,}9997$ \\
GRPO few-shot                  & $0{,}3601$ & $0{,}0162$ & $0{,}3325$ & $0{,}9991$ \\
SFT$\to$GRPO zero-shot         & $\mathbf{0{,}5274}$ & $\mathbf{0{,}0283}$ & $\mathbf{0{,}4599}$ & $1{,}0000$ \\
SFT$\to$GRPO few-shot          & $0{,}5039$ & $0{,}0240$ & $0{,}4429$ & $1{,}0000$ \\
GRPO senza Trie, few-shot      & $0{,}3550$ & $0{,}0165$ & $0{,}3339$ & $0{,}9916$ \\
GRPO senza Trie, zero-shot     & $0{,}2898$ & $0{,}0000$ & $0{,}0000$ & $0{,}0000$ \\
SFT con massa ammessa          & $0{,}4819$ & $0{,}0178$ & $0{,}3709$ & $0{,}9981$ \\
\bottomrule
\end{tabular}
\caption{PHOENIX-2014T, test completo.}
\label{tab:risultati-phoenix}
\end{table}

Su PHOENIX il margine c'è, e nessun sistema si avvicina al tetto: l'exact match
resta sotto il $3\%$ ovunque. Qui, a differenza di ASLG, il GRPO dopo l'SFT
aggiunge qualcosa di visibile: ROUGE-L da $0{,}4854$ a $0{,}5274$ e non-copy da
$0{,}3812$ a $0{,}4599$, dello stesso ordine dei migliori valori in
letteratura~\cite{ouargani2023t2g} (con implementazioni di ROUGE diverse, per
cui il confronto è solo indicativo). Il GRPO dal modello base impara (ROUGE-L
da $0{,}0961$ a $0{,}2030$ in zero-shot, da $0{,}3103$ a $0{,}3601$ in
few-shot) ma resta sotto l'SFT. La massa ammessa, che su ASLG era la cella
migliore, qui vale quanto l'SFT semplice.

La cella GRPO senza Trie zero-shot ha ROUGE-L $0{,}29$ con validity zero: il
modello scrive parole tedesche minuscole separate da virgole
(\texttt{regen, schnee, nacht, nach, nordosten, schauer}), che ROUGE-L,
insensibile al maiuscolo, premia per accidente (sezione~\ref{sec:difetto-rouge}).

\subsection{CoNLL-2003}

\begin{table}[h]
\centering
\small
\begin{tabular}{lrrrr}
\toprule
Cella & Exact match & Gloss-F1 & ROUGE-L & Validity \\
\midrule
Base zero-shot                 & $0{,}0453$ & $0{,}0893$ & $0{,}2708$ & $0{,}9693$ \\
Base few-shot                  & $0{,}0662$ & $0{,}1429$ & $0{,}3805$ & $0{,}9739$ \\
Base zero-shot senza Trie      & $0{,}0011$ & $0{,}0089$ & $0{,}2676$ & $0{,}0079$ \\
SFT                            & $0{,}7432$ & $0{,}8330$ & $0{,}9003$ & $0{,}9998$ \\
GRPO zero-shot                 & $0{,}0228$ & $0{,}1964$ & $0{,}3526$ & $0{,}9979$ \\
GRPO few-shot                  & $0{,}3315$ & $0{,}3885$ & $0{,}5483$ & $0{,}9979$ \\
SFT$\to$GRPO zero-shot         & $\mathbf{0{,}7807}$ & $\mathbf{0{,}8494}$ & $\mathbf{0{,}9128}$ & $0{,}9999$ \\
SFT$\to$GRPO few-shot          & $0{,}7289$ & $0{,}8230$ & $0{,}8859$ & $0{,}9995$ \\
GRPO senza Trie, few-shot      & $0{,}3473$ & $0{,}4205$ & $0{,}5906$ & $0{,}7261$ \\
GRPO senza Trie, zero-shot     & $0{,}0611$ & $0{,}0142$ & $0{,}4548$ & $0{,}1430$ \\
SFT con massa ammessa          & $0{,}7469$ & $0{,}8359$ & $0{,}9039$ & $1{,}0000$ \\
\bottomrule
\end{tabular}
\caption{CoNLL-2003, $2.000$ prompt del test, \emph{prima} della correzione
della sezione~\ref{sec:span-unici}. Gloss-F1 è il micro-F1 sui token
\texttt{TIPO:entità}, cioè l'F1 a livello di entità.}
\label{tab:risultati-conll}
\end{table}

La gerarchia è netta: l'SFT porta l'exact match da $0{,}05$ a $0{,}74$, il GRPO
dopo l'SFT aggiunge quasi quattro punti in zero-shot ($0{,}7807$), il GRPO dal
modello base resta lontano. La cella GRPO zero-shot è quella dell'hack della
sezione~\ref{sec:span-unici}: è l'unica in cui il GRPO \emph{peggiora} l'exact
match rispetto al punto di partenza. In few-shot, dove gli esempi mostrano il
formato e una sola etichetta per entità, l'hack è molto più raro ($8\%$ delle
uscite) e l'exact match sale a $0{,}33$.

\subsection{WOS-46985}

\begin{table}[h]
\centering
\small
\begin{tabular}{lrrr}
\toprule
Cella & Exact match & Gloss-F1 & Validity \\
\midrule
Base zero-shot                 & $0{,}1039$ & $0{,}2876$ & $1{,}0000$ \\
Base zero-shot senza Trie      & $0{,}0000$ & $0{,}0146$ & $0{,}0000$ \\
SFT                            & $0{,}7622$ & $0{,}8269$ & $1{,}0000$ \\
GRPO zero-shot                 & $0{,}1688$ & $0{,}3702$ & $1{,}0000$ \\
SFT$\to$GRPO zero-shot         & $\mathbf{0{,}7869}$ & $\mathbf{0{,}8475}$ & $1{,}0000$ \\
GRPO senza Trie                & $0{,}0000$ & $0{,}1467$ & $0{,}0001$ \\
SFT con massa ammessa          & $0{,}7572$ & $0{,}8255$ & $1{,}0000$ \\
\bottomrule
\end{tabular}
\caption{WOS-46985, $2.000$ prompt del test, solo zero-shot.}
\label{tab:risultati-wos}
\end{table}

Con il Trie a sequenze il modello base è sempre valido ma sceglie quasi a caso
fra poche classi: \texttt{ECE System\_identification} e \texttt{Psychology
Social\_cognition} coprono da sole un terzo delle sue uscite. Il GRPO migliora di
sei punti ($0{,}1039 \to 0{,}1688$) senza uscire da questo regime; l'SFT porta
l'exact match a $0{,}76$, e il GRPO dopo l'SFT aggiunge due punti e mezzo.
Senza Trie il modello non produce mai una coppia valida: scrive
\texttt{Medical Psychology, Medical Psychology}, con dominio e area presi
dalla stessa lista e una virgola.

\section{Che cosa dicono i tre dataset insieme}
\label{sec:osservazioni-altri}

\subsection{Senza Trie e senza esempi, il GRPO non trova il formato}

È il risultato più diretto per la domanda del progetto. Senza Trie, la validity
delle celle GRPO dipende tutta dalla presenza di esempi nel prompt:

\begin{center}
\small
\begin{tabular}{lrr}
\toprule
GRPO senza Trie & few-shot & zero-shot \\
\midrule
ASLG-PC12     & $0{,}99$ & --- \\
PHOENIX-2014T & $0{,}99$ & $0{,}00$ \\
CoNLL-2003    & $0{,}73$ & $0{,}14$ \\
WOS-46985     & ---      & $0{,}0001$ \\
\bottomrule
\end{tabular}
\end{center}

Il meccanismo si legge nei log di addestramento. Su PHOENIX zero-shot la
componente di formato della reward resta fra $-1$ e $-0{,}99$ per tutti i
$5.000$ passi: nessun candidato ha mai il formato giusto, quindi quasi nessun
gruppo contiene un confronto fra uscita valida e non valida, e il formato non
riceve gradiente. Su CoNLL e WOS
qualche candidato valido compare per caso e la componente sale (da $-0{,}89$ a
$-0{,}31$ e da $-0{,}98$ a $-0{,}12$), ma non abbastanza da rendere valide le
uscite in valutazione. Il GRPO, dunque, \textbf{rinforza una struttura che il
modello già produce qualche volta; non la scopre}. Gli esempi nel prompt, il
Trie o l'SFT sono i supporti che la rendono presente fin dall'inizio, e senza
almeno uno dei tre la struttura non arriva.

\subsection{Con un vincolo attivo, il GRPO cerca il varco}

L'hack di CoNLL (sezione~\ref{sec:span-unici}) è il rovescio della stessa
medaglia: quando la grammatica lascia un'uscita ammessa che la reward premia più
del dovuto, il GRPO la trova, e la validity non lo segnala perché l'uscita è
formalmente corretta. Un vincolo rigido va quindi scritto pensando a che cosa
\emph{non} deve permettere, non solo a che cosa deve garantire.

\subsection{SFT$\to$GRPO è la cella migliore ovunque, ma il guadagno è piccolo}

Su tutti e tre i dataset la cella migliore è l'SFT seguito dal GRPO in
zero-shot, con un margine di due-quattro punti di exact match sull'SFT. Come
su ASLG, il segnale di apprendimento si esaurisce: alla fine dell'addestramento
la frazione di gruppi a varianza nulla è l'$85\%$ su CoNLL e il $77\%$ su WOS
(su PHOENIX, dove la policy è più lontana dal tetto, solo il $7\%$). La massa
ammessa, migliore cella su ASLG-PC12, sugli altri tre dataset vale quanto l'SFT
semplice.

\begin{table}[h]
\centering
\small
\begin{tabular}{lrrrr}
\toprule
 & \multicolumn{2}{c}{Reward media} & \multicolumn{2}{c}{Gruppi a varianza nulla} \\
Cella & inizio & fine & inizio & fine \\
\midrule
PHOENIX GRPO zero-shot        & $-0{,}60$ & $-0{,}39$ & $0\%$  & $0\%$  \\
PHOENIX SFT$\to$GRPO zero-shot & $-0{,}06$ & $0{,}10$ & $0\%$  & $7\%$  \\
CoNLL GRPO zero-shot          & $-0{,}49$ & $-0{,}28$ & $18\%$ & $18\%$ \\
CoNLL SFT$\to$GRPO zero-shot  & $0{,}76$ & $0{,}82$  & $72\%$ & $85\%$ \\
WOS GRPO zero-shot            & $-0{,}38$ & $-0{,}21$ & $34\%$ & $36\%$ \\
WOS SFT$\to$GRPO zero-shot    & $0{,}56$ & $0{,}68$  & $58\%$ & $77\%$ \\
\bottomrule
\end{tabular}
\caption{Medie sui primi e sugli ultimi $50$ passi registrati di ogni
addestramento GRPO (dati di CoNLL precedenti alla correzione).}
\label{tab:dinamica-altri}
\end{table}
